% ============================================================================
% DISCUSSION SECTION FOR DDI KNOWLEDGE GRAPH AND POLYPHARMACY RISK ASSESSMENT
% Publication Format - Structured Discussion
% ============================================================================

\documentclass[11pt,a4paper]{article}
\usepackage[utf8]{inputenc}
\usepackage{amsmath,amssymb}
\usepackage{booktabs}
\usepackage{geometry}
\usepackage{graphicx}
\usepackage{hyperref}
\usepackage{float}
\usepackage{setspace}
\geometry{margin=1in}
\onehalfspacing

\hypersetup{
    colorlinks=true,
    linkcolor=blue,
    filecolor=magenta,
    urlcolor=cyan,
    citecolor=blue
}

\begin{document}

\section*{Discussion}

Cardiovascular disease remains the leading cause of mortality worldwide, with polypharmacy affecting over 40\% of elderly cardiac patients \cite{january2019afib}. We focused on this domain because it exemplifies the polypharmacy challenge: narrow therapeutic indices (Warfarin, Digoxin), extensive CYP450-mediated interactions (Amiodarone), and life-threatening consequences of both under- and over-treatment \cite{steinberg2017off}. Notably, three of the top five hub drugs in our knowledge graph are cardiovascular agents, empirically confirming this domain's disproportionate interaction burden. The 759,774-interaction dataset represents one of the largest structured DDI repositories assembled, yet even this scale likely underrepresents pharmacodynamic interactions that are inherently difficult to catalog \cite{tatonetti2012data}.

Key methodological strengths include fact-based linkage ensuring specificity without spurious relationships \cite{han2022systematic}, proper train/test validation avoiding data leakage \cite{kapoor2023leakage}, and interpretable composite scoring addressing ``black box'' criticisms of clinical AI \cite{rudin2019stop}. However, the weak Cohen's $\kappa$ (0.096) reflects fundamental challenges in text-based severity classification---substantial keyword overlap across severity levels creates an inherent ceiling that more sophisticated NLP may partially but not fully overcome \cite{brown2020language}. The heuristically-defined PRI thresholds require prospective validation against patient outcomes before clinical deployment.

Our work differs from prior efforts in prioritizing specificity over sensitivity. While similarity-based DDI prediction achieves $>$90\% accuracy \cite{ryu2018deep, deng2020multimodal, chen2021transformer}, such methods introduce false-positives that exacerbate alert fatigue---a critical limitation given that 49--96\% of DDI alerts are currently overridden \cite{van2006override}. The network-centric PRI differs from drug-counting indices like the Anticholinergic Cognitive Burden scale \cite{campbell2009drug} by capturing systemic risk propagation through centrality metrics, aligning with network medicine principles \cite{barabasi2011network}. Critically, our recalibration reveals that zero-shot transformer classification assigns 99.9\% of interactions as Major or Contraindicated---grossly inconsistent with clinical experience and a likely contributor to alert fatigue \cite{phansalkar2013drug}.

For clinicians, the Dabigatran recommendation for Warfarin substitution aligns with contemporary AF guidelines \cite{january2019afib, joglar2023afib}, providing computational validation for evidence-based practice. The recalibrated severity distribution (80.1\% Moderate vs. DrugBank's 0.003\%) may enable more discriminating alert strategies. For policymakers, simple rule-based classification outperforming complex transformers (66.4\% vs. 22.1\%) has implications for democratizing access in resource-limited settings, while fact-based provenance addresses regulatory transparency requirements \cite{fda2021ai}.

Future work should prioritize prospective PRI validation against adverse events and hospitalizations \cite{kashyap2019clinical}; enhancement through fine-tuned LLMs or patient-specific factors \cite{singhal2023large, johnson2011clinical}; domain-specific validation in oncology and psychiatry; and human factors research on optimal alert thresholds \cite{kawamoto2005improving}. This study establishes a foundation for evidence-based polypharmacy optimization, demonstrating that systematic computational approaches can support clinicians in managing multi-drug therapy complexity. All recommendations require verification against established clinical evidence before decision-making.

% ============================================================================
% REFERENCES
% ============================================================================

\bibliographystyle{vancouver}
\begin{thebibliography}{99}

\bibitem{barabasi2004network}
Barab\'{a}si AL, Oltvai ZN.
Network biology: understanding the cell's functional organization.
\textit{Nat Rev Genet}. 2004;5(2):101-113. doi:10.1038/nrg1272

\bibitem{barabasi2011network}
Barab\'{a}si AL, Gulbahce N, Loscalzo J.
Network medicine: a network-based approach to human disease.
\textit{Nat Rev Genet}. 2011;12(1):56-68. doi:10.1038/nrg2918

\bibitem{brown2020language}
Brown TB, Mann B, Ryder N, et al.
Language models are few-shot learners.
\textit{Adv Neural Inf Process Syst}. 2020;33:1877-1901. arXiv:2005.14165

\bibitem{campbell2009drug}
Campbell N, Boustani M, Limbil T, et al.
The cognitive impact of anticholinergics: a clinical review.
\textit{Clin Interv Aging}. 2009;4:225-233. doi:10.2147/cia.s5358

\bibitem{chen2021transformer}
Chen X, Liu X, Wu J.
Drug-drug interaction prediction with transformer-based models.
\textit{Brief Bioinform}. 2021;22(6):bbab203. doi:10.1093/bib/bbab203

\bibitem{deng2020multimodal}
Deng Y, Xu X, Qiu Y, et al.
A multimodal deep learning framework for predicting drug-drug interaction events.
\textit{Bioinformatics}. 2020;36(15):4316-4322. doi:10.1093/bioinformatics/btaa501

\bibitem{fda2021ai}
U.S. Food and Drug Administration.
Artificial Intelligence/Machine Learning (AI/ML)-Based Software as a Medical Device (SaMD) Action Plan.
FDA; January 2021. Available at: https://www.fda.gov/medical-devices/software-medical-device-samd/artificial-intelligence-and-machine-learning-software-medical-device

\bibitem{han2022systematic}
Han K, Cao P, Wang Y, et al.
A review on drug-drug interaction prediction using machine learning.
\textit{Drug Discov Today}. 2022;27(4):1095-1107. doi:10.1016/j.drudis.2021.12.008

\bibitem{hilmer2007drug}
Hilmer SN, Mager DE, Simonsick EM, et al.
A drug burden index to define the functional burden of medications in older people.
\textit{Arch Intern Med}. 2007;167(8):781-787. doi:10.1001/archinte.167.8.781

\bibitem{january2019afib}
January CT, Wann LS, Calkins H, et al.
2019 AHA/ACC/HRS Focused Update of the 2014 AHA/ACC/HRS Guideline for the Management of Patients With Atrial Fibrillation.
\textit{J Am Coll Cardiol}. 2019;74(1):104-132. doi:10.1016/j.jacc.2019.01.011

\bibitem{joglar2023afib}
Joglar JA, Chung MK, Armbruster AL, et al.
2023 ACC/AHA/ACCP/HRS Guideline for the Diagnosis and Management of Atrial Fibrillation.
\textit{J Am Coll Cardiol}. 2024;83(1):109-279. doi:10.1016/j.jacc.2023.08.017

\bibitem{johnson2011clinical}
Johnson JA, Gong L, Whirl-Carrillo M, et al.
Clinical Pharmacogenetics Implementation Consortium Guidelines for CYP2C9 and VKORC1 genotypes and warfarin dosing.
\textit{Clin Pharmacol Ther}. 2011;90(4):625-629. doi:10.1038/clpt.2011.185

\bibitem{kapoor2023leakage}
Kapoor S, Narayanan A.
Leakage and the reproducibility crisis in machine-learning-based science.
\textit{Patterns}. 2023;4(9):100804. doi:10.1016/j.patter.2023.100804

\bibitem{kashyap2019clinical}
Kashyap R, Dutt T, Engoren M.
Machine learning methodologies in clinical prediction and decision support.
\textit{Crit Care Med}. 2019;47(7):946-949. doi:10.1097/CCM.0000000000003767

\bibitem{kawamoto2005improving}
Kawamoto K, Houlihan CA, Balas EA, Lobach DF.
Improving clinical practice using clinical decision support systems: a systematic review of trials to identify features critical to success.
\textit{BMJ}. 2005;330(7494):765. doi:10.1136/bmj.38398.500764.8F

\bibitem{kuhn2016sider}
Kuhn M, Letunic I, Jensen LJ, Bork P.
The SIDER database of drugs and side effects.
\textit{Nucleic Acids Res}. 2016;44(D1):D1075-D1079. doi:10.1093/nar/gkv1075

\bibitem{lin2020kgnn}
Lin X, Quan Z, Wang ZJ, et al.
KGNN: Knowledge Graph Neural Network for Drug-Drug Interaction Prediction.
\textit{Proc Int Joint Conf Artif Intell}. 2020;2739-2745. doi:10.24963/ijcai.2020/380

\bibitem{marsh2016multiple}
Marsh K, Lanitis T, Neasham D, et al.
Assessing the value of healthcare interventions using multi-criteria decision analysis: a review of the literature.
\textit{Pharmacoeconomics}. 2014;32(4):345-365. doi:10.1007/s40273-014-0135-0

\bibitem{phansalkar2013drug}
Phansalkar S, Desai AA, Bell D, et al.
High-priority drug-drug interactions for use in electronic health records.
\textit{J Am Med Inform Assoc}. 2012;19(5):735-743. doi:10.1136/amiajnl-2011-000612

\bibitem{rudin2019stop}
Rudin C.
Stop explaining black box machine learning models for high stakes decisions and use interpretable models instead.
\textit{Nat Mach Intell}. 2019;1(5):206-215. doi:10.1038/s42256-019-0048-x

\bibitem{ryu2018deep}
Ryu JY, Kim HU, Lee SY.
Deep learning improves prediction of drug-drug and drug-food interactions.
\textit{Proc Natl Acad Sci USA}. 2018;115(18):E4304-E4311. doi:10.1073/pnas.1803294115

\bibitem{singhal2023large}
Singhal K, Azizi S, Tu T, et al.
Large language models encode clinical knowledge.
\textit{Nature}. 2023;620(7972):172-180. doi:10.1038/s41586-023-06291-2

\bibitem{steinberg2017off}
Steinberg BA, Shrader P, Thomas L, et al.
Off-label dosing of non-vitamin K antagonist oral anticoagulants and adverse outcomes: the ORBIT-AF II registry.
\textit{J Am Coll Cardiol}. 2016;68(24):2597-2604. doi:10.1016/j.jacc.2016.09.966

\bibitem{tatonetti2012data}
Tatonetti NP, Ye PP, Daneshjou R, Altman RB.
Data-driven prediction of drug effects and interactions.
\textit{Sci Transl Med}. 2012;4(125):125ra31. doi:10.1126/scitranslmed.3003377

\bibitem{van2006override}
van der Sijs H, Aarts J, Vulto A, Berg M.
Overriding of drug safety alerts in computerized physician order entry.
\textit{J Am Med Inform Assoc}. 2006;13(2):138-147. doi:10.1197/jamia.M1809

\bibitem{wishart2018drugbank}
Wishart DS, Feunang YD, Guo AC, et al.
DrugBank 5.0: a major update to the DrugBank database for 2018.
\textit{Nucleic Acids Res}. 2018;46(D1):D1074-D1082. doi:10.1093/nar/gkx1037

\bibitem{wright2019reduced}
Wright A, Ash JS, Erickson JL, et al.
A qualitative study of the activities performed by people involved in clinical decision support: recommended practices for success.
\textit{J Am Med Inform Assoc}. 2014;21(3):464-472. doi:10.1136/amiajnl-2013-001771

\bibitem{xiong2022ddinter}
Xiong G, Yang Z, Yi J, et al.
DDInter: an online drug-drug interaction database towards improving clinical decision-making and patient safety.
\textit{Nucleic Acids Res}. 2022;50(D1):D1200-D1207. doi:10.1093/nar/gkab880

\bibitem{zitnik2018modeling}
Zitnik M, Agrawal M, Leskovec J.
Modeling polypharmacy side effects with graph convolutional networks.
\textit{Bioinformatics}. 2018;34(13):i457-i466. doi:10.1093/bioinformatics/bty294

\end{thebibliography}

\end{document}
